\SourceLeaf{009}
\Article{2}
Was bisher dargelegt wurde, wird in den Büchern über Algebra ausreichend bewiesen und verstößt nirgends gegen die geometrische Strenge. Doch scheinen die Analysten allzu voreilig und ohne vorherigen stichhaltigen Beweis den Satz angenommen zu haben, auf dem fast die gesamte Lehre von den Gleichungen beruht: \textit{Jede Funktion wie $X$ lässt sich stets in $m$ lineare Faktoren zerlegen}, oder, was damit völlig übereinstimmt, \textit{jede Gleichung vom Grad $m$ hat tatsächlich $m$ Wurzeln}. Da schon bei quadratischen Gleichungen sehr häufig Fälle auftreten, die diesem Satz widersprechen, waren die Algebraiker gezwungen, um auch diese Fälle unter ihn zu bringen, eine imaginäre Größe zu erfinden, deren Quadrat $-1$ ist. Darauf erkannten sie: Lässt man Größen der Form $a+b\sqrt{-1}$ ebenso zu wie reelle, so gilt der Satz nicht nur für quadratische, sondern auch für kubische und biquadratische Gleichungen. Daraus durfte man jedoch keineswegs schließen, dass sich nach Zulassung der Größen der Form $a+b\sqrt{-1}$ jede Gleichung fünften oder höheren Grades erfüllen lässt oder, wie man gewöhnlich sagt (obwohl ich diese schlüpfrige Formulierung nicht gutheißen möchte), dass die Wurzeln jeder Gleichung auf die Form $a+b\sqrt{-1}$ gebracht werden können. Der Sache nach unterscheidet sich dieser Satz in nichts von dem im Titel dieser Schrift ausgesprochenen; einen neuen strengen Beweis dafür zu geben, ist das Ziel der vorliegenden Dissertation.

Seit die Analysten erkannt haben, dass es unendlich viele Gleichungen gibt, die überhaupt keine Wurzel hätten, wenn nicht Größen der Form $a+b\sqrt{-1}$ zugelassen würden, hat man solche erdachten Größen als eine besondere Art von Größen betrachtet, sie zur Unterscheidung von den \textit{reellen} als \textit{imaginäre} bezeichnet und in die gesamte Analysis eingeführt. Mit welchem Recht? Das erörtere ich hier nicht. --- Ich werde meinen Beweis ohne jede Hilfe imaginärer Größen durchführen, obwohl es auch mir erlaubt wäre, dieselbe Freiheit in Anspruch zu nehmen, deren sich alle neueren Analysten bedient haben.

\Article{3}
Was in den meisten elementaren Büchern als Beweis unseres Satzes angeführt wird, ist zwar so gehaltlos und von geometrischer Strenge so weit entfernt, dass es kaum der Erwähnung wert ist; damit aber nichts fehlt, will ich es dennoch kurz berühren. `Um zu zeigen, dass jede Gleichung
\[
x^m+Ax^{m-1}+Bx^{m-2}+\text{usw.}+M=0
\]
\SourceLeaf{010}
oder $X=0$ tatsächlich $m$ Wurzeln hat, unternimmt man den Beweis, dass sich $X$ in $m$ lineare Faktoren zerlegen lässt. Dazu nimmt man $m$ lineare Faktoren $x-\alpha$, $x-\beta$, $x-\gamma$ usw. an, wobei $\alpha$, $\beta$, $\gamma$ usw. noch unbekannt sind, und setzt ihr Produkt der Funktion $X$ gleich. Durch Vergleich der Koeffizienten erhält man dann $m$ Gleichungen, aus denen sich, wie man sagt, die Unbekannten $\alpha$, $\beta$, $\gamma$ usw. bestimmen lassen, da auch ihre Anzahl $m$ ist. Man könne nämlich $m-1$ Unbekannte eliminieren, sodass eine Gleichung übrig bleibe, die nur noch diejenige Unbekannte enthalte, die man beliebig auswählt.' Von den übrigen Einwänden gegen eine solche Schlussweise will ich schweigen und nur fragen: Woher können wir sicher wissen, dass diese letzte Gleichung tatsächlich irgendeine Wurzel hat? Warum könnte es nicht sein, dass weder diese letzte noch die vorgelegte Gleichung durch irgendeine Größe im gesamten Bereich der reellen und imaginären Größen erfüllt wird? --- Im Übrigen werden Sachkundige leicht erkennen, dass diese letzte Gleichung notwendig mit der vorgelegten \textit{völlig identisch} ist, sofern die \mbox{Rechnung richtig durchgeführt wurde}: Nach Elimination der Unbekannten $\beta$, $\gamma$ usw. muss nämlich die Gleichung
\[
\alpha^m+A\alpha^{m-1}+B\alpha^{m-2}+\text{usw.}+M=0
\]
entstehen. Mehr braucht über diese Schlussweise nicht gesagt zu werden.

Einige Autoren, die die Schwäche dieser Methode bemerkt zu haben scheinen, nehmen als \textit{Axiom} an, jede Gleichung habe tatsächlich Wurzeln, wenn nicht mögliche, dann unmögliche. Was sie unter möglichen und unmöglichen Größen verstehen wollen, scheinen sie nicht hinreichend deutlich erläutert zu haben. Sollen die möglichen Größen dasselbe wie die reellen, die unmöglichen dasselbe wie die imaginären bezeichnen, so darf jenes Axiom keineswegs zugestanden werden, sondern bedarf notwendig eines Beweises. Doch scheinen die Ausdrücke nicht in diesem Sinn gemeint zu sein; vielmehr scheint der Sinn des Axioms dieser zu sein: `Obwohl wir noch nicht sicher sind, dass es notwendig $m$ reelle oder imaginäre Größen gibt, die eine gegebene Gleichung vom Grad $m$ erfüllen, wollen wir dies dennoch eine Zeit lang voraussetzen; sollte es nämlich geschehen, dass sich so viele reelle und imaginäre Größen nicht finden lassen, so bleibt uns gewiss der Ausweg zu sagen, die übrigen seien unmöglich.' Zieht jemand diese Ausdrucksweise der einfachen Aussage vor, die Gleichung werde in diesem Fall nicht so viele Wurzeln haben, so habe ich nichts dagegen. Benutzt er dann aber diese unmöglichen Wurzeln, als wären sie etwas Wirkliches, und sagt zum Beispiel, die Summe aller Wurzeln der Gleichung $x^m+Ax^{m-1}+\text{usw.}=0$ sei $=-A$, auch wenn sich unmögliche darunter befinden (was eigentlich bedeutet: \textit{auch wenn einige fehlen}), so kann ich dies keineswegs gutheißen. Denn unmögliche Wurzeln\SourceLeaf{011}, in diesem Sinn genommen, sind dennoch Wurzeln; dann darf jenes Axiom auf keinen Fall ohne Beweis zugestanden werden, und wäre es abwegig zu fragen, ob es nicht Gleichungen geben könnte, die nicht einmal unmögliche Wurzeln haben?\HistoricalNote{1}

\Article{4}
Bevor ich die Beweise anderer Mathematiker für unseren Satz durchgehe und darlege, was mir an jedem von ihnen zu beanstanden scheint, bemerke ich: Es genügt, nur zu zeigen, dass sich jede Gleichung eines beliebigen Grades
\[
x^m+Ax^{m-1}+Bx^{m-2}+\text{usw.}+M=0
\]
oder $X=0$ (wobei die Koeffizienten $A,B$ usw. als reell vorausgesetzt werden) auf wenigstens eine Weise durch einen Wert von $x$ der Form $a+b\sqrt{-1}$ erfüllen lässt. Bekanntlich ist $X$ dann durch den reellen quadratischen Faktor \mbox{$xx-2ax+aa+bb$} teilbar, wenn $b$ nicht $=0$ ist, und durch den reellen linearen Faktor $x-a$, wenn $b=0$ ist. In beiden Fällen ist der Quotient reell und von niedrigerem Grad als $X$; und da dieser auf dieselbe Weise einen reellen Faktor ersten oder zweiten Grades haben muss\SourceLeaf{012}, ist klar, dass die Fortsetzung dieses Verfahrens die Funktion $X$ schließlich in reelle lineare oder quadratische Faktoren zerlegt oder, wenn man für jeden reellen quadratischen Faktor lieber zwei imaginäre lineare verwendet, in $m$ lineare Faktoren.
